\documentclass[a4paper, serif, 10pt]{moderncv}

%% moderncv
% color and style
\moderncvcolor{black}
\moderncvstyle{banking}

%% page layout/margins
\usepackage[top=2.5cm, bottom=2.5cm, left=1.5cm, right=1.5cm]{geometry}

\nopagenumbers{}

%% Babel config for Italian language
% ref:
% - https://latex3.github.io/babel/guides/locale-italian.html
\usepackage[italian]{babel}

%% fontspec
\usepackage{fontspec}

\IfFontExistsTF{Linux Libertine}{
  \setmainfont[Ligatures={TeX, Common}, Numbers=OldStyle]{Linux Libertine}
}{}
\IfFontExistsTF{Linux Libertine O}{
  \setmainfont[Ligatures={TeX, Common}, Numbers=OldStyle]{Linux Libertine O}
}{}

%% CTeX
% CJK support, used to show CJK characters in the resume
%
% - heading=false: disable ctex heading
% - punct=kaiming: use kaiming punctuations styles for CJK
% - scheme=plain: use plain scheme, do not override \normalsize font size
% - space=auto: space settings for CJK characters
%
% CJK fonts are resolved per locale language, e.g. Noto Serif/Sans CJK JP
% for Japanese resumes, with fallback to the next region if not installed
%
% ref:
% - http://ctan.mirrorcatalogs.com/language/chinese/ctex/ctex.pdf
\usepackage[UTF8, heading=false, punct=kaiming, scheme=plain, space=auto]{ctex}

\IfFontExistsTF{Noto Serif CJK SC}{
  \setCJKmainfont{Noto Serif CJK SC}
}{}
\IfFontExistsTF{Noto Sans CJK SC}{
  \setCJKsansfont{Noto Sans CJK SC}
}{}

%% line spacing
\usepackage{setspace}
\setstretch{1.125}

%% URL styling - use normal text instead of monospace
\urlstyle{same}

%% Patch \httplink and \httpslink to handle full URLs
% The original moderncv macros always prepend http:// or https:// to URLs.
% This causes issues when the URL already has a protocol (e.g., https://example.com)
% resulting in malformed URLs like https://https://example.com.
% This patch checks if the URL already contains "://" and uses it directly if so.
% Note: We use \str_set:Nx (x-expansion) to expand macros like \@homepage first.
\ExplSyntaxOn
\str_new:N \g_yamlresume_url_str

\RenewDocumentCommand{\httpslink}{O{}m}{
  \str_set:Nx \g_yamlresume_url_str {#2}
  \str_if_in:NnTF \g_yamlresume_url_str {://}
    { \url{#2} }
    { \url{https://#2} }
}

\RenewDocumentCommand{\httplink}{O{}m}{
  \str_set:Nx \g_yamlresume_url_str {#2}
  \str_if_in:NnTF \g_yamlresume_url_str {://}
    { \url{#2} }
    { \url{http://#2} }
}
\ExplSyntaxOff

\name{Marco Bianchi}{}
\title{Ingegnere DevOps specializzato in cloud e automazione}
\phone[mobile]{+39 333 123 4567}
\email{marco.bianchi@example.com}
\homepage{https://www.marcobianchi.dev}

\address{Via Alessandro Manzoni 15, Milano, Lombardia, Italia, 20121}{}{}

\extrainfo{{\small \faGithub}\ \href{https://github.com/marcobianchi-devops}{@marcobianchi-devops} {} {} {} • {} {} {} 
{\small \faLinkedin}\ \href{https://www.linkedin.com/in/marco-bianchi-devops}{@marco-bianchi-devops} {} {} {} • {} {} {} 
{\small \faStackOverflow}\ \href{https://stackoverflow.com/users/12345678/marco-bianchi}{@marco-bianchi}}

\begin{document}

\maketitle

\section{Informazioni personali}

\cvline{}{Ingegnere DevOps con oltre 7 anni di esperienza nella progettazione e gestione di infrastrutture cloud scalabili e sicure.
%
\begin{itemize}
\item Esperto in automazione dei processi di build, deploy e monitoraggio su AWS e Kubernetes
\item Forte orientamento alla cultura DevOps, al miglioramento continuo e all'affidabilità dei sistemi
\item Competenze consolidate in Infrastructure as Code, CI/CD e osservabilità
\item Abituato a collaborare in team Agile e a mentorare sviluppatori su pratiche SRE
\end{itemize}}
\section{Formazione}

\cventry{set 2010–lug 2012}
        {Laurea magistrale, Ingegneria Informatica, Punteggio: 108/110}
        {Politecnico di Milano}
        {\url{https://www.polimi.it}}
        {}
        {\begin{itemize}
\item Specializzazione in sistemi distribuiti e architetture cloud
\item Tesi sperimentale su orchestrazione di container e scalabilità automatica
\item Partecipazione a progetti di ricerca su affidabilità dei sistemi
\end{itemize}
\textbf{Corsi}: Sistemi Distribuiti, Cloud Computing, Sicurezza Informatica, Ingegneria del Software, Basi di Dati, Reti di Calcolatori, Architetture dei Sistemi Software}

\cventry{set 2007–lug 2010}
        {Laurea triennale, Ingegneria Informatica, Punteggio: 105/110}
        {Università degli Studi di Bologna}
        {\url{https://www.unibo.it}}
        {}
        {\begin{itemize}
\item Solida preparazione di base in informatica e ingegneria del software
\item Progetti di gruppo su sviluppo di applicazioni web e basi di dati
\end{itemize}
\textbf{Corsi}: Programmazione, Algoritmi e Strutture Dati, Sistemi Operativi, Basi di Dati, Reti di Calcolatori, Ingegneria del Software}
\section{Esperienza professionale}

\cventry{mar 2019–Presente}
        {DevOps Engineer}
        {TechNova S.p.A.}
        {\url{https://www.technova.it}}
        {}
        {\begin{itemize}
\item Progettazione e gestione di infrastrutture cloud su AWS con Terraform e Ansible
\item Implementazione di pipeline CI/CD con Jenkins, GitLab CI e Argo CD per oltre 40 microservizi
\item Migrazione di workload on-premise a Kubernetes (EKS) riducendo i costi operativi del 30\%
\item Definizione di metriche SLO/SLI e dashboard di monitoraggio con Prometheus e Grafana
\item Automazione di processi di sicurezza e compliance con Vault e OPA
\end{itemize}
\textbf{Parole chiave}: AWS, Terraform, Kubernetes, CI/CD, Prometheus}

\cventry{giu 2016–feb 2019}
        {Site Reliability Engineer}
        {CloudSphere Italia}
        {\url{https://www.cloudsphere.it}}
        {}
        {\begin{itemize}
\item Gestione di piattaforme Kubernetes multi-tenant su GCP e Azure
\item Riduzione del MTTR attraverso alerting proattivo e runbook automatizzati
\item Sviluppo di operatori Kubernetes in Go per il self-service dei team di sviluppo
\item Implementazione di policy di autoscaling e capacity planning
\end{itemize}
\textbf{Parole chiave}: GCP, Azure, Kubernetes, Go, SRE}

\cventry{set 2012–mag 2016}
        {Junior System Administrator}
        {FinData Solutions S.r.l.}
        {\url{https://www.findata.it}}
        {}
        {\begin{itemize}
\item Amministrazione di sistemi Linux e Windows in ambienti ad alta disponibilità
\item Automazione di attività ripetitive con script Bash e Python
\item Gestione di backup, disaster recovery e monitoraggio con Nagios e Zabbix
\item Supporto alla migrazione verso infrastrutture virtualizzate VMware
\end{itemize}
\textbf{Parole chiave}: Linux, Bash, Python, VMware, Nagios}
\section{Lingue}

\cvline{Italiano}{Madrelingua o bilingue \hfill \textbf{Parole chiave}: Madrelingua}
\cvline{Inglese}{Competenza professionale completa \hfill \textbf{Parole chiave}: IELTS 8.0, TOEFL 115}
\cvline{Tedesco}{Competenza professionale limitata \hfill \textbf{Parole chiave}: B1}
\section{Competenze}

\cvline{Cloud \& Infrastructure}{Esperto \hfill \textbf{Parole chiave}: AWS, GCP, Azure, Terraform, Ansible}
\cvline{Container \& Orchestration}{Esperto \hfill \textbf{Parole chiave}: Kubernetes, Docker, Helm, Argo CD, EKS}
\cvline{CI/CD \& Automation}{Avanzato \hfill \textbf{Parole chiave}: Jenkins, GitLab CI, GitHub Actions, Bash, Python}
\cvline{Monitoring \& Observability}{Avanzato \hfill \textbf{Parole chiave}: Prometheus, Grafana, ELK Stack, Datadog, OpenTelemetry}
\cvline{Security \& Compliance}{Intermedio \hfill \textbf{Parole chiave}: Vault, OPA, IAM, SOC 2, GDPR}
\section{Premi}

\cventry{nov 2021}
        {TechNova S.p.A.}
        {Premio Innovazione Tecnologica}
        {}
        {}
        {Riconoscimento per il progetto di migrazione cloud che ha ridotto i costi infrastrutturali del 30\% e migliorato l'affidabilità dei servizi.}
\section{Certificazioni}

\cventry{mar 2022}
        {Amazon Web Services}
        {AWS Certified DevOps Engineer – Professional}
        {\url{https://aws.amazon.com/certification/certified-devops-engineer-professional/}}
        {}
        {}

\cventry{lug 2020}
        {Cloud Native Computing Foundation}
        {Certified Kubernetes Administrator (CKA)}
        {\url{https://www.cncf.io/certification/cka/}}
        {}
        {}

\cventry{feb 2021}
        {HashiCorp}
        {HashiCorp Certified: Terraform Associate}
        {\url{https://www.hashicorp.com/certification/terraform-associate}}
        {}
        {}
\section{Pubblicazioni}

\cventry{apr 2020}
        {Pratiche SRE per piattaforme Kubernetes multi-tenant}
        {DevOps Italia Magazine}
        {\url{https://www.devopsitalia.it/articoli/sre-kubernetes}}
        {}
        {\begin{itemize}
\item Analisi delle strategie di osservabilità e gestione degli incidenti
\item Case study su riduzione del MTTR e automazione dei runbook
\end{itemize}}
\section{Progetti}

\cventry{gen 2021–dic 2022}
        {Portale interno per il provisioning automatizzato di ambienti cloud}
        {Piattaforma di self-service per team di sviluppo}
        {\url{https://github.com/marcobianchi-devops/devportal}}
        {}
        {\begin{itemize}
\item Progettazione di un portale basato su Backstage per l'accesso self-service a risorse cloud
\item Integrazione con Terraform e Argo CD per il deploy automatico degli ambienti
\item Riduzione del time-to-market dei team di sviluppo del 40\%
\end{itemize}
\textbf{Parole chiave}: Backstage, Terraform, Argo CD, Self-service}
\section{Interessi}

\cvline{Open Source}{Kubernetes, Terraform, Prometheus}
\cvline{Sport}{Corsa, Ciclismo, Nuoto}
\section{Volontariato}

\cventry{gen 2018–Presente}
        {Mentor e relatore}
        {DevOps Milano}
        {\url{https://www.devopsmilano.it}}
        {}
        {\begin{itemize}
\item Mentoring per professionisti junior su pratiche DevOps e cloud
\item Organizzazione di meetup tecnici e workshop su Kubernetes e Infrastructure as Code
\item Relatore su temi di osservabilità e SRE
\end{itemize}}
    
\end{document}